\section{模式四：多轮对话评测 --- 模拟真实交互}

\subsection{多轮评测的核心挑战}

真实场景中，用户往往与 Agent 进行多轮交互。例如先问``帮我安排明天的会议''，然后根据 Agent 的回复进一步说``改到下午 3 点''。然而，多轮评测面临一个根本性问题：

\begin{warningbox}{硬编码多轮输入的脆弱性}
如果天真地将多轮对话的用户输入硬编码为一个固定序列，当 Agent 在第一轮的回复偏离预期时，后续的硬编码输入将\textbf{与上下文脱节}，导致整个测试变得无意义。例如，如果 Agent 没有按预期询问``你想安排什么时间？''，而是直接建议了一个时间，那么硬编码的``下午 3 点''回复就毫无意义了。
\end{warningbox}

\subsection{条件分支策略}

LangChain 团队采用的解决方案是在 Pytest/Vitest 测试中引入\textbf{条件逻辑}：

\begin{enumerate}
  \item 运行第一轮，检查 Agent 输出
  \item 如果输出符合预期，继续运行下一轮
  \item 如果输出不符合预期，\textbf{提前终止测试}（early fail）
\end{enumerate}

这种方式的优势在于：
\begin{itemize}
  \item 不需要穷举 Agent 的所有可能分支
  \item 每一轮都有独立的检查点（checkpoint）
  \item 失败时可以精确定位到哪一轮出了问题
\end{itemize}

\begin{knowledgebox}{隔离测试特定轮次}
如果需要单独测试第二轮或第三轮的行为，可以直接用适当的初始状态（initial state）启动测试，跳过前面的轮次。这在 LangGraph 中特别方便，因为 Agent 的状态可以被序列化和恢复。这种技巧让多轮测试的编写和维护成本大幅降低。
\end{knowledgebox}

\subsection{多轮评测的设计建议}

根据 LangChain 的经验，多轮评测的设计应遵循以下原则：

\begin{table}[H]
\centering
\caption{多轮评测设计原则}
\begin{tabular}{lp{9cm}}
\toprule
\textbf{原则} & \textbf{说明} \\
\midrule
最少轮次 & 每个测试用例只包含验证目标行为所需的最少轮数 \\
条件推进 & 每一轮结束后验证输出，只在预期路径上推进 \\
早期失败 & 不符合预期时立即终止，避免浪费资源在无效状态上运行 \\
状态快照 & 支持从任意中间状态恢复，便于隔离测试特定轮次 \\
\bottomrule
\end{tabular}
\end{table}

\subsection{本章小结}

多轮对话评测模拟了真实的用户交互场景，但面临输入与上下文脱节的挑战。通过条件分支策略和状态快照机制，可以在保持测试有效性的同时控制复杂度。关键是\textbf{不要试图穷举所有对话路径}，而是选择最关键的场景进行有针对性的测试。


